\section{Precision, Recall and Stability by Language}\label{app:tables}
\begin{table}[!ht]
\centering
\caption{Boundary precision (\%).}
\label{tab:precision}
\small
\setlength{\tabcolsep}{4pt}
\begin{tabular}{lccccc}
\toprule
System & English & Spanish & Danish & Tamil & Turkish \\
\midrule
BPE, free marker & 31.6\,{\scriptsize$\pm$1.2} & 13.7\,{\scriptsize$\pm$0.5} & 17.9\,{\scriptsize$\pm$0.6} & 3.2\,{\scriptsize$\pm$1.7} & 24.7\,{\scriptsize$\pm$1.4} \\
POS-weighted BPE & 33.6\,{\scriptsize$\pm$1.1} & 13.8\,{\scriptsize$\pm$0.6} & 18.1\,{\scriptsize$\pm$0.3} & 4.5\,{\scriptsize$\pm$1.7} & 24.6\,{\scriptsize$\pm$1.3} \\
\quad transitions only & 33.6\,{\scriptsize$\pm$1.2} & 13.8\,{\scriptsize$\pm$0.6} & 18.1\,{\scriptsize$\pm$0.3} & 4.6\,{\scriptsize$\pm$1.7} & 24.6\,{\scriptsize$\pm$1.3} \\
\quad penalty only & 32.7\,{\scriptsize$\pm$0.9} & 13.7\,{\scriptsize$\pm$0.5} & 18.0\,{\scriptsize$\pm$0.7} & 3.2\,{\scriptsize$\pm$1.7} & 24.6\,{\scriptsize$\pm$1.5} \\
\addlinespace[2pt]
Content-word-weighted BPE & 34.3\,{\scriptsize$\pm$0.2} & 12.1\,{\scriptsize$\pm$0.2} & 20.1\,{\scriptsize$\pm$0.5} & 16.8\,{\scriptsize$\pm$0.9} & 24.2\,{\scriptsize$\pm$0.3} \\
Lemma2Word + BPE & \textbf{45.3\,{\scriptsize$\pm$0.1}} & \textbf{25.6\,{\scriptsize$\pm$0.1}} & \textbf{31.7\,{\scriptsize$\pm$0.2}} & \textbf{38.5\,{\scriptsize$\pm$0.4}} & 34.4\,{\scriptsize$\pm$0.1} \\
BPE, word-initial marker & 22.0\,{\scriptsize$\pm$0.3} & 5.3\,{\scriptsize$\pm$0.4} & 10.2\,{\scriptsize$\pm$0.1} & 6.4\,{\scriptsize$\pm$0.3} & 31.0\,{\scriptsize$\pm$0.3} \\
Unigram LM & 45.0\,{\scriptsize$\pm$0.5} & 21.9\,{\scriptsize$\pm$0.0} & 27.2\,{\scriptsize$\pm$0.4} & 36.9\,{\scriptsize$\pm$0.8} & \textbf{35.0\,{\scriptsize$\pm$0.3}} \\
\bottomrule
\end{tabular}
\end{table}

\begin{table}[!ht]
\centering
\caption{Boundary recall (\%).}
\label{tab:recall}
\small
\setlength{\tabcolsep}{4pt}
\begin{tabular}{lccccc}
\toprule
System & English & Spanish & Danish & Tamil & Turkish \\
\midrule
BPE, free marker & 54.4\,{\scriptsize$\pm$2.6} & 29.4\,{\scriptsize$\pm$1.5} & 37.9\,{\scriptsize$\pm$1.5} & 7.4\,{\scriptsize$\pm$3.7} & 60.2\,{\scriptsize$\pm$2.4} \\
POS-weighted BPE & 61.3\,{\scriptsize$\pm$2.5} & 29.8\,{\scriptsize$\pm$1.5} & 37.9\,{\scriptsize$\pm$1.3} & 9.5\,{\scriptsize$\pm$3.3} & 60.7\,{\scriptsize$\pm$2.1} \\
\quad transitions only & 61.3\,{\scriptsize$\pm$2.5} & 29.8\,{\scriptsize$\pm$1.5} & 37.9\,{\scriptsize$\pm$1.3} & 9.6\,{\scriptsize$\pm$3.3} & 60.7\,{\scriptsize$\pm$2.1} \\
\quad penalty only & 58.7\,{\scriptsize$\pm$2.0} & 29.4\,{\scriptsize$\pm$1.5} & 37.9\,{\scriptsize$\pm$1.6} & 7.4\,{\scriptsize$\pm$3.7} & 60.1\,{\scriptsize$\pm$2.5} \\
\addlinespace[2pt]
Content-word-weighted BPE & 60.8\,{\scriptsize$\pm$0.3} & 25.9\,{\scriptsize$\pm$0.5} & 44.8\,{\scriptsize$\pm$1.1} & 32.7\,{\scriptsize$\pm$1.6} & 57.8\,{\scriptsize$\pm$0.4} \\
Lemma2Word + BPE & \textbf{88.2\,{\scriptsize$\pm$0.1}} & \textbf{62.2\,{\scriptsize$\pm$0.3}} & \textbf{70.8\,{\scriptsize$\pm$0.5}} & \textbf{84.5\,{\scriptsize$\pm$0.5}} & \textbf{83.3\,{\scriptsize$\pm$0.3}} \\
BPE, word-initial marker & 30.3\,{\scriptsize$\pm$0.4} & 9.1\,{\scriptsize$\pm$0.7} & 18.4\,{\scriptsize$\pm$0.2} & 9.0\,{\scriptsize$\pm$0.4} & 59.6\,{\scriptsize$\pm$0.7} \\
Unigram LM & 68.6\,{\scriptsize$\pm$0.4} & 42.9\,{\scriptsize$\pm$0.2} & 49.3\,{\scriptsize$\pm$0.6} & 51.5\,{\scriptsize$\pm$0.9} & 69.3\,{\scriptsize$\pm$0.7} \\
\bottomrule
\end{tabular}
\end{table}

\begin{table}[!ht]
\centering
\caption{Share of held-out tokens that span a word boundary (\%).}
\label{tab:cross}
\small
\setlength{\tabcolsep}{4pt}
\begin{tabular}{lccccc}
\toprule
System & English & Spanish & Danish & Tamil & Turkish \\
\midrule
BPE, free marker & 15.2\,{\scriptsize$\pm$0.1} & 18.3\,{\scriptsize$\pm$0.0} & 13.1\,{\scriptsize$\pm$0.1} & \textbf{22.0\,{\scriptsize$\pm$1.5}} & 4.5\,{\scriptsize$\pm$0.1} \\
POS-weighted BPE & 14.2\,{\scriptsize$\pm$0.1} & 18.4\,{\scriptsize$\pm$0.0} & 13.2\,{\scriptsize$\pm$0.2} & 15.6\,{\scriptsize$\pm$2.4} & 4.4\,{\scriptsize$\pm$0.1} \\
\quad transitions only & 14.2\,{\scriptsize$\pm$0.1} & 18.4\,{\scriptsize$\pm$0.0} & 13.2\,{\scriptsize$\pm$0.2} & 15.6\,{\scriptsize$\pm$2.4} & 4.4\,{\scriptsize$\pm$0.1} \\
\quad penalty only & 14.4\,{\scriptsize$\pm$0.0} & 18.3\,{\scriptsize$\pm$0.0} & 13.2\,{\scriptsize$\pm$0.1} & 21.9\,{\scriptsize$\pm$1.5} & 4.5\,{\scriptsize$\pm$0.1} \\
\addlinespace[2pt]
Content-word-weighted BPE & 10.9\,{\scriptsize$\pm$0.2} & 14.5\,{\scriptsize$\pm$0.1} & 9.7\,{\scriptsize$\pm$0.2} & 2.2\,{\scriptsize$\pm$0.3} & 2.5\,{\scriptsize$\pm$0.1} \\
Lemma2Word + BPE & \textbf{15.7\,{\scriptsize$\pm$0.0}} & \textbf{19.2\,{\scriptsize$\pm$0.0}} & \textbf{15.4\,{\scriptsize$\pm$0.1}} & 12.8\,{\scriptsize$\pm$0.2} & \textbf{6.6\,{\scriptsize$\pm$0.1}} \\
BPE, word-initial marker & 0.0\,{\scriptsize$\pm$0.0} & 0.0\,{\scriptsize$\pm$0.0} & 0.0\,{\scriptsize$\pm$0.0} & 0.0\,{\scriptsize$\pm$0.0} & 0.0\,{\scriptsize$\pm$0.0} \\
Unigram LM & 0.0\,{\scriptsize$\pm$0.0} & 0.0\,{\scriptsize$\pm$0.0} & 0.0\,{\scriptsize$\pm$0.0} & 0.0\,{\scriptsize$\pm$0.0} & 0.0\,{\scriptsize$\pm$0.0} \\
\bottomrule
\end{tabular}
\end{table}

\begin{table}[!ht]
\centering
\caption{Vocabulary reuse: Jaccard overlap (\%) between the vocabularies learned from different training samples (three seed pairs).}
\label{tab:jaccard}
\small
\setlength{\tabcolsep}{4pt}
\begin{tabular}{lccccc}
\toprule
System & English & Spanish & Danish & Tamil & Turkish \\
\midrule
BPE, free marker & 76.8\,{\scriptsize$\pm$0.8} & 75.8\,{\scriptsize$\pm$1.8} & 76.8\,{\scriptsize$\pm$0.4} & 63.8\,{\scriptsize$\pm$4.8} & 74.6\,{\scriptsize$\pm$1.9} \\
POS-weighted BPE & 77.8\,{\scriptsize$\pm$0.9} & 76.3\,{\scriptsize$\pm$2.0} & 76.7\,{\scriptsize$\pm$0.6} & 61.9\,{\scriptsize$\pm$5.4} & 67.1\,{\scriptsize$\pm$7.4} \\
\quad transitions only & 77.8\,{\scriptsize$\pm$0.9} & 76.2\,{\scriptsize$\pm$2.0} & 76.6\,{\scriptsize$\pm$0.5} & 61.8\,{\scriptsize$\pm$5.4} & 67.2\,{\scriptsize$\pm$7.4} \\
\quad penalty only & 77.1\,{\scriptsize$\pm$0.7} & 75.8\,{\scriptsize$\pm$1.8} & 76.3\,{\scriptsize$\pm$0.4} & 64.1\,{\scriptsize$\pm$4.8} & 75.1\,{\scriptsize$\pm$2.3} \\
\addlinespace[2pt]
Content-word-weighted BPE & 74.8\,{\scriptsize$\pm$3.1} & \textbf{77.2\,{\scriptsize$\pm$0.3}} & 74.3\,{\scriptsize$\pm$2.4} & 71.5\,{\scriptsize$\pm$2.5} & 77.0\,{\scriptsize$\pm$1.0} \\
Lemma2Word + BPE & 76.9\,{\scriptsize$\pm$0.7} & 75.6\,{\scriptsize$\pm$1.5} & \textbf{77.8\,{\scriptsize$\pm$0.6}} & 74.8\,{\scriptsize$\pm$0.5} & 74.1\,{\scriptsize$\pm$3.1} \\
BPE, word-initial marker & \textbf{78.5\,{\scriptsize$\pm$0.4}} & 74.4\,{\scriptsize$\pm$3.2} & 76.1\,{\scriptsize$\pm$0.9} & \textbf{76.9\,{\scriptsize$\pm$0.2}} & \textbf{78.5\,{\scriptsize$\pm$1.1}} \\
Unigram LM & 69.3\,{\scriptsize$\pm$1.1} & 66.3\,{\scriptsize$\pm$0.7} & 62.7\,{\scriptsize$\pm$0.2} & 62.9\,{\scriptsize$\pm$0.6} & 69.1\,{\scriptsize$\pm$0.4} \\
\bottomrule
\end{tabular}
\end{table}

\begin{table}[!ht]
\centering
\caption{Segmentation agreement: boundary F1 (\%) between the held-out segmentations of tokenizers trained on different samples (three seed pairs).}
\label{tab:agreement}
\small
\setlength{\tabcolsep}{4pt}
\begin{tabular}{lccccc}
\toprule
System & English & Spanish & Danish & Tamil & Turkish \\
\midrule
BPE, free marker & 95.2\,{\scriptsize$\pm$0.1} & 95.1\,{\scriptsize$\pm$0.8} & 95.7\,{\scriptsize$\pm$0.2} & 88.6\,{\scriptsize$\pm$3.4} & 94.6\,{\scriptsize$\pm$1.2} \\
POS-weighted BPE & 95.8\,{\scriptsize$\pm$0.1} & 95.2\,{\scriptsize$\pm$0.9} & 95.5\,{\scriptsize$\pm$0.3} & 87.5\,{\scriptsize$\pm$3.3} & 92.5\,{\scriptsize$\pm$2.7} \\
\quad transitions only & 95.8\,{\scriptsize$\pm$0.1} & 95.3\,{\scriptsize$\pm$0.9} & 95.5\,{\scriptsize$\pm$0.3} & 87.6\,{\scriptsize$\pm$3.4} & 92.5\,{\scriptsize$\pm$2.7} \\
\quad penalty only & 95.6\,{\scriptsize$\pm$0.2} & 95.1\,{\scriptsize$\pm$0.8} & 95.5\,{\scriptsize$\pm$0.1} & 88.7\,{\scriptsize$\pm$3.5} & 94.7\,{\scriptsize$\pm$1.3} \\
\addlinespace[2pt]
Content-word-weighted BPE & 95.0\,{\scriptsize$\pm$1.1} & \textbf{95.9\,{\scriptsize$\pm$0.1}} & 95.5\,{\scriptsize$\pm$0.6} & 94.0\,{\scriptsize$\pm$0.9} & \textbf{96.0\,{\scriptsize$\pm$0.2}} \\
Lemma2Word + BPE & 95.2\,{\scriptsize$\pm$0.3} & 94.9\,{\scriptsize$\pm$0.6} & 95.8\,{\scriptsize$\pm$0.2} & 94.7\,{\scriptsize$\pm$0.2} & 95.6\,{\scriptsize$\pm$0.7} \\
BPE, word-initial marker & \textbf{96.9\,{\scriptsize$\pm$0.1}} & 95.5\,{\scriptsize$\pm$1.2} & \textbf{95.9\,{\scriptsize$\pm$0.1}} & \textbf{95.4\,{\scriptsize$\pm$0.2}} & 95.8\,{\scriptsize$\pm$0.6} \\
Unigram LM & 95.0\,{\scriptsize$\pm$0.1} & 94.1\,{\scriptsize$\pm$0.2} & 93.5\,{\scriptsize$\pm$0.2} & 93.5\,{\scriptsize$\pm$0.1} & 94.3\,{\scriptsize$\pm$0.2} \\
\bottomrule
\end{tabular}
\end{table}


\section{Example Segmentations}\label{app:examples}
\begin{table}[!ht]
\centering
\caption{Example segmentations of evaluation words absent from the training text (seed 0). Words are the first, middle and last of the alphabetically sorted unseen forms without capitals of 8 to 14 characters with at least one gold boundary; the word-initial-marker BPE is omitted for space. Vertical bars mark boundaries; the gold column shows the lemma-derived boundaries.}
\label{tab:examples}
\scriptsize
\setlength{\tabcolsep}{2.5pt}
\begin{tabular}{llllll}
\toprule
Lang. & Gold & BPE, free marker & Content-weighted & Lemma2Word & Unigram \\
\midrule
En & \texttt{abduct|ed} & \texttt{ab|duct|ed} & \texttt{ab|duct|ed} & \texttt{ab|duc|t|ed} & \texttt{a|b|du|c|ted} \\
En & \texttt{interrupt|ing} & \texttt{inter|rupt|ing} & \texttt{inter|rupt|ing} & \texttt{inter|rup|t|ing} & \texttt{interrupt|ing} \\
En & \texttt{youngster|s} & \texttt{young|st|ers} & \texttt{young|st|ers} & \texttt{young|ster|s} & \texttt{young|ster|s} \\
\addlinespace[2pt]
Sp & \texttt{abandonado|s} & \texttt{abandon|ados} & \texttt{abandon|ados} & \texttt{abandon|a|dos} & \texttt{abandonado|s} \\
Sp & \texttt{galeote|s} & \texttt{gal|e|ot|es} & \texttt{gal|e|ot|es} & \texttt{gal|e|ot|es} & \texttt{gal|eo|tes} \\
Sp & \texttt{ímprobo|s} & \texttt{í|m|prob|os} & \texttt{í|m|prob|os} & \texttt{í|m|prob|os} & \texttt{í|mp|ro|bo|s} \\
\addlinespace[2pt]
Da & \texttt{abonnent|erne} & \texttt{ab|on|n|ent|erne} & \texttt{abon|n|ent|ern|e} & \texttt{ab|on|n|ent|erne} & \texttt{ab|on|n|en|t|erne} \\
Da & \texttt{landshold|ets} & \texttt{lands|hol|det|s} & \texttt{lands|hold|et|s} & \texttt{lands|hold|et|s} & \texttt{lands|hold|ets} \\
Da & \texttt{østmagt|ernes} & \texttt{øst|magt|ern|es} & \texttt{øst|magt|ern|es} & \texttt{øst|magt|erne|s} & \texttt{øst|magt|ernes} \\
\addlinespace[2pt]
Ta & \tamil{gஎர்மனி|யில்} & \tamil{g|எ|ர்மன|ிய|ில்} & \tamil{g|எ|ர்மன|ிய|ில்} & \tamil{g|எ|ர்ம|னி|யில்} & \tamil{g|எ|ர்|ம|னி|யில்} \\
Ta & \tamil{தற்கொலை|ப்} & \tamil{தற்க|ொலைப|்} & \tamil{தற்கொ|லை|ப்} & \tamil{தற்க|ொ|ல|ை|ப்} & \tamil{தற்கொலை|ப்} \\
Ta & \tamil{ஹென்றி|யின்} & \tamil{ஹ|ென்ற|ிய|ின்} & \tamil{ஹெ|ன்ற|ிய|ின்} & \tamil{ஹெ|ன்ற|ி|ய|ின்} & \tamil{ஹென்றி|யின்} \\
\addlinespace[2pt]
Tu & \texttt{abajur|lara} & \texttt{ab|aj|ur|lar|a} & \texttt{ab|aj|ur|lar|a} & \texttt{aba|j|ur|lar|a} & \texttt{a|ba|j|ur|lara} \\
Tu & \texttt{korkala|mıştı} & \texttt{kor|kal|amış|tı} & \texttt{kor|kal|amış|t|ı} & \texttt{kor|kal|amış|tı} & \texttt{kor|ka|la|mıştı} \\
Tu & \texttt{şıpıdık|larının} & \texttt{ş|ıp|ı|dık|ların|ın} & \texttt{ş|ıp|ıd|ık|larının} & \texttt{ş|ıp|ı|dık|ların|ın} & \texttt{ş|ıp|ı|dıkları|nın} \\
\addlinespace[2pt]
\bottomrule
\end{tabular}
\end{table}


\section{Prompt of the Segmentation Study}\label{app:prompt}
The language model received this text, with the language name and the batch of words filled in, one word per line.
\begin{quote}\small\ttfamily\raggedright
Segment each \{lang\} word below into its morphemes: prefixes, the stem, and suffixes (inflectional and derivational).\\
Rules: keep every character of the word exactly as given and in order; insert a plus sign ({+}) at each morpheme boundary; a word with a single morpheme is returned unchanged; do not add explanations.\\
Answer with one line per word in the format \ \ word: seg{+}men{+}ted \ \ and nothing else.\par\medskip
\{words\}
\end{quote}

\section{Reproduction}\label{app:repro}
All code, the data-preparation scripts and the per-run metrics are in the repository.
The commands below regenerate every number, table and figure.
\begin{verbatim}
python experiments/fetch_corpus.py          # Wikipedia paragraphs
bash   experiments/tag_all.sh               # Stanza tags and lemmas
bash   experiments/run_all.sh               # 5 languages x 8 methods x 3 seeds
bash   experiments/run_sweep.sh             # 2,000 and 32,000 types, seed 0
bash   experiments/after_grid.sh            # language-model probes
python experiments/llm_segment.py           # LLM segmentation study
python experiments/aggregate.py             # numbers.tex, tables, figures
\end{verbatim}
The gold boundaries come from the Hugging Face dataset \path{catherinearnett/morphscore} and the text from \path{wikimedia/wikipedia}.
